% ============================================================================
% LANGENTWURF LE-014: E-Mail verfassen
% Profil: S (Senioren 65+) | Tier: 2 (Standard)
% ============================================================================

\documentclass[11pt,a4paper]{article}
\usepackage[utf8]{inputenc}
\usepackage[T1]{fontenc}
\usepackage[ngerman]{babel}
\usepackage{geometry}
\usepackage{fancyhdr}
\usepackage{lastpage}
\usepackage{xcolor}
\usepackage{tcolorbox}
\usepackage{enumitem}
\usepackage{longtable}
\usepackage{hyperref}
\usepackage{amssymb}

\geometry{left=2.5cm, right=2.5cm, top=2.5cm, bottom=2.5cm}
\definecolor{fach}{HTML}{b35c00}
\definecolor{gray}{HTML}{666666}
\definecolor{successgreen}{HTML}{2a7a4b}
\definecolor{warnred}{HTML}{c62828}

\tcbuselibrary{skins,breakable}
\newtcolorbox{scorebox}{ colback=white, colframe=fach, fonttitle=\bfseries, title=Didaktik-Score, breakable }
\newtcolorbox{tippbox}[1][]{ colback=successgreen!5, colframe=successgreen, fonttitle=\bfseries, title=#1, breakable }

\pagestyle{fancy}
\fancyhf{}
\fancyhead[L]{\small\textcolor{gray}{Digitale Grundlagen -- 014 E-Mail verfassen}}
\fancyhead[R]{\small\textcolor{gray}{LE Tier 2 / Profil S}}
\fancyfoot[C]{\small\thepage{} / \pageref{LastPage}}
\renewcommand{\headrulewidth}{0.4pt}

\begin{document}

\begin{titlepage}
    \centering
    \vspace*{2cm}
    {\Large\textcolor{gray}{Digitale Grundlagen -- Senioren 65+}}\\[2cm]
    {\Huge\textcolor{fach}{\textbf{Langentwurf}}}\\[0.5cm]
    {\large Tier 2 | Profil S}\\[2cm]
    {\LARGE\textbf{014 -- E-Mail verfassen \& senden}}\\[0.5cm]
    {\large Modul C: Internet \& Kommunikation}\\[2cm]
    \begin{tabular}{rl}
        \textbf{Zeitrahmen:} & 60--75 Minuten \\
        \textbf{Vorkenntnisse:} & E-Mail verstehen (013) \\
    \end{tabular}
    \vfill
    {\normalsize Stand: \today}
\end{titlepage}

\tableofcontents
\newpage

\section{Bedingungsanalyse}

\subsection{Ausgangssituation}
\begin{itemize}
    \item E-Mails lesen können viele, schreiben ist schwieriger
    \item Unsicherheit bei Anrede und Grußformel
    \item Angst, ``etwas falsch zu machen''
    \item Anhänge hinzufügen ist unklar
    \item CC und BCC sind unbekannte Abkürzungen
\end{itemize}

\subsection{Typische Probleme}
\begin{itemize}
    \item ``Wie fange ich an?''
    \item ``Wie hänge ich ein Foto an?''
    \item ``An wen geht die Mail eigentlich?''
    \item ``Kann ich die Mail zurückholen, wenn ich sie abgeschickt habe?''
\end{itemize}

\section{Sachanalyse}

\subsection{Kernkonzepte}
\begin{enumerate}
    \item \textbf{An (To):} Der Hauptempfänger -- die Person, für die die Mail bestimmt ist
    \item \textbf{CC (Kopie):} Weitere Empfänger, die die Mail zur Info bekommen
    \item \textbf{Betreff:} Kurze Überschrift, worum es geht
    \item \textbf{Anrede:} Wie beim Brief -- ``Sehr geehrte...'', ``Liebe/Lieber...''
    \item \textbf{Signatur:} Automatischer Abschluss mit Name und Kontaktdaten
    \item \textbf{Anhang:} Dateien, die du mitschickst
\end{enumerate}

\subsection{Alltagsanalogien}
\begin{tabular}{|l|p{8cm}|}
\hline
\textbf{Begriff} & \textbf{Analogie} \\
\hline
An (To) & Der Name auf dem Briefumschlag \\
CC & Kopie an weitere Personen (sichtbar für alle) \\
BCC & Kopie an weitere Personen (geheim, andere sehen es nicht) \\
Betreff & Die Überschrift auf einer Postkarte \\
Anhang & Etwas, das du dem Brief beilegst \\
Senden & Den Brief in den Briefkasten werfen \\
\hline
\end{tabular}

\section{Didaktische Analyse}

\subsection{Stellung in der Reihe}
Aufbauend auf 013 (E-Mail verstehen). Aktive Kommunikationsfähigkeit.

\subsection{Didaktische Reduktion}
\textbf{Ausgeklammert:} BCC im Detail, Prioritäten, Lesebestätigungen, HTML-Formatierung.

\textbf{Fokus:} Eine einfache E-Mail schreiben und senden, Foto anhängen.

\section{Lernziele}

\begin{tabular}{|c|p{7cm}|c|c|}
\hline
\textbf{Nr.} & \textbf{Die Teilnehmer können...} & \textbf{Stufe} & \textbf{Niveau} \\
\hline
FZ1 & eine neue E-Mail beginnen & Anwenden & Basis \\
FZ2 & Empfänger, Betreff und Text eingeben & Anwenden & Basis \\
FZ3 & eine passende Anrede und Grußformel wählen & Anwenden & Basis \\
FZ4 & ein Foto als Anhang hinzufügen & Anwenden & Basis \\
FZ5 & auf eine E-Mail antworten oder weiterleiten & Anwenden & Vertiefung \\
\hline
\end{tabular}

\section{Methodische Analyse}

\subsection{E-Mail-Aufbau}
\begin{tippbox}[So schreiben Sie eine E-Mail]
\begin{enumerate}
    \item \textbf{An:} E-Mail-Adresse des Empfängers eingeben
    \item \textbf{Betreff:} Kurz und klar, z.B. ``Terminanfrage Zahnarzt''
    \item \textbf{Anrede:} ``Sehr geehrte Frau/Herr...'' oder ``Liebe/Lieber...''
    \item \textbf{Text:} Das, was Sie mitteilen möchten
    \item \textbf{Grußformel:} ``Mit freundlichen Grüßen'' oder ``Liebe Grüße''
    \item \textbf{Name:} Ihr Name am Ende
\end{enumerate}
\end{tippbox}

\section{Verlaufsplanung}

\begin{longtable}{|p{1cm}|p{2cm}|p{5cm}|p{3cm}|p{2cm}|}
\hline
\textbf{Zeit} & \textbf{Phase} & \textbf{Inhalt} & \textbf{TN-Aktivität} & \textbf{Material} \\
\hline
0--10 & Einstieg & Wann schreiben Sie Briefe? E-Mail als digitaler Brief & Austauschen & -- \\
\hline
10--25 & Neue Mail & Neue-Mail-Symbol, An, Betreff erklären, Demo & Mitmachen & S.1 \\
\hline
25--40 & Text schreiben & Anrede, Text, Grußformel, Tastatur-Tipps & Üben & S.1 \\
\hline
40--55 & Anhang & Foto aus Fotos-App anhängen, Größenhinweis & Selbst versuchen & S.2 \\
\hline
55--70 & Antworten & Auf Mail antworten, weiterleiten, Übungen & Bearbeiten & S.2-3 \\
\hline
\end{longtable}

\section{Lösungen}

\textbf{Neue E-Mail schreiben (iPhone):}
\begin{itemize}
    \item Mail-App öffnen
    \item Unten rechts: Quadrat mit Stift (Neue Mail)
    \item Oben ``An:'' eintippen oder aus Kontakten wählen
    \item ``Betreff:'' kurze Überschrift
    \item Darunter: Text schreiben
    \item Oben rechts: ``Senden'' (blauer Pfeil)
\end{itemize}

\textbf{Foto anhängen:}
\begin{itemize}
    \item Im Textfeld tippen und halten
    \item ``Foto oder Video einfügen'' wählen
    \item ODER: Büroklammer-Symbol / < -Symbol antippen
    \item Foto aus Mediathek wählen
\end{itemize}

\textbf{Anrede-Beispiele:}
\begin{itemize}
    \item \textbf{Formell:} ``Sehr geehrte Frau Müller,'' / ``Sehr geehrter Herr Schmidt,''
    \item \textbf{Halbformell:} ``Guten Tag Frau Müller,''
    \item \textbf{Informell:} ``Liebe Maria,'' / ``Hallo Peter,''
\end{itemize}

\section{Didaktik-Score}

\begin{scorebox}
\begin{tabular}{|c|l|c|c|p{3.5cm}|}
\hline
\# & Dimension & Gew. & Score & Notiz \\
\hline
1 & Differenzierung & 15\% & 8/10 & Basis/Vertiefung klar \\
2 & Taxonomie & 10\% & 9/10 & Schwerpunkt Anwenden \\
3 & Alltagsrelevanz & 10\% & 10/10 & Wichtige Kommunikation \\
4 & Aktivierung & 10\% & 10/10 & Eigene Mail schreiben \\
5 & Scaffolding & 10\% & 10/10 & Schritt-für-Schritt-Box \\
6 & Feedback & 10\% & 8/10 & Checkliste, Beispiele \\
7 & Sprachliche Klarheit & 10\% & 9/10 & Einfach, Brief-Analogie \\
8 & Motivation & 10\% & 9/10 & Kommunikation motiviert \\
9 & Barrierefreiheit & 10\% & 9/10 & Große Darstellung \\
10 & Praktikabilität & 5\% & 9/10 & 70 Min angemessen \\
\hline
\multicolumn{3}{|r|}{\textbf{GESAMT:}} & \textbf{9.15/10} & \textcolor{successgreen}{\textbf{FREIGABE}} \\
\hline
\end{tabular}
\end{scorebox}

\end{document}
